\chapter*{Preface}
The prime numbers look random. They thin out, but irregularly; twin primes keep appearing, and long gaps too; no
formula produces them. Yet in bulk they obey precise laws. Their number up to $x$ is asymptotic to
$\frac x{\log x}$; they are shared equally among the residue classes modulo $q$ that can contain them; the number of
prime factors of a typical integer follows a Gaussian law. This volume of \emph{Order and Chance}, which closes the
first part of the series, proves these laws, and it is where the earlier volumes meet: the integers of Volume~2, the
characters of Volume~3, the probability of Volume~4 and the analysis of Volume~5 are used together.

Chapters~1 and~2 are elementary estimates: Mertens's theorems on sums over primes, and the average size of
arithmetic functions. Chapters~3 and~4 build Dirichlet's $L$-functions and prove that every progression $a+nq$ with
$\gcd(a,q)=1$ contains infinitely many primes. Chapter~5 proves the prime number theorem with Newman's short analytic
argument; Chapters~6 and~7 study the zeros of the zeta function, give the theorem an error term and derive
Riemann's explicit formula, in which the primes appear as a superposition of waves, one for each zero. Chapters~8--12
are sieves, probabilistic number theory and the Erd\H{o}s--Kac theorem, random models of the primes, exponential sums,
and a first taste of modular forms. The book stops at the classical explicit formula, Selberg's sieve, the Erd\H{o}s--Kac theorem and the shape of Cramér's model. Bombieri--Vinogradov, the circle method and a positive proportion of zeros on the critical line are Volume~14.

Throughout, $\log$ is the natural logarithm, as is usual in number theory, and $p$ always denotes a prime.

Every theorem, proposition and lemma is proved; results quoted only for their interest appear as remarks, with a
reference to where they are proved. Exercises are graded \textbf{$\star$} for a warm-up, \textbf{$\star\star$} for a good
afternoon and \textbf{$\star\star\star$} for gems, and every one has a hint and a complete solution at the end of its
chapter, beginning with \emph{How to find it} before the \emph{Write-up}. ``Pause and think'' boxes ask you to guess
before you read on. The only fact used without proof is the Fubini--Tonelli theorem for general integrands; each use is
labelled \emph{Admitted (proved in Volume~7)}.

The series asks whether pure order and pure chance can be told apart, and why the primes look random yet obey laws. This floor is the span: the prime number theorem, the explicit formula, and the Gaussian law of $\omega(n)$. The next floor takes up the words the span used without defining them: a limit, a length, a room. Maier's theorem is marked and not yet proved.
